% Slideaway — Seminar/Invited Talk Template (45–60 min)
% Theme: Madrid/dolphin — softer blue tone for longer departmental talks
% Structural scaffolding only. Style/color theory → design-foundations.md
% Engine workflow → engine-beamer.md

\documentclass[aspectratio=169,11pt]{beamer}

\usetheme{Madrid}
\usecolortheme{dolphin}

\setbeamertemplate{navigation symbols}{}
\setbeamertemplate{footline}[frame number]

% Section dividers — pause between major segments
\AtBeginSection[]{
  \begin{frame}
    \vfill
    \centering
    \begin{beamercolorbox}[sep=8pt,center,shadow=true,rounded=true]{title}
      \usebeamerfont{title}\insertsectionhead\par
    \end{beamercolorbox}
    \vfill
  \end{frame}
}

% ── Packages ──────────────────────────────────────────────────────────────
\usepackage{booktabs}
\usepackage{graphicx}
\usepackage{amsmath}
\usepackage{hyperref}
\usepackage{xcolor}
\usepackage{algorithm}
\usepackage{algorithmic}
\usepackage{listings}
\usepackage{multirow}
\usepackage[backend=biber,style=authoryear]{biblatex}
% \addbibresource{references.bib}

% Code listing defaults
\lstset{
  basicstyle=\ttfamily\scriptsize,
  keywordstyle=\color{blue},
  commentstyle=\color{gray},
  breaklines=true,
  frame=single,
  numbers=left,
  numberstyle=\tiny\color{gray},
}

% ── Title Metadata ────────────────────────────────────────────────────────
\title[Short Title]{Full Seminar Title}
\subtitle{Department Seminar Series — Spring 2026}
\author[A.~Author]{Author Name}
\institute{Department of X \\ University of Y}
\date{\today}

% ══════════════════════════════════════════════════════════════════════════
\begin{document}

\begin{frame}
  \titlepage
\end{frame}

\begin{frame}{Outline}
  \tableofcontents
\end{frame}

% ── Motivation and Background (5–7 slides, ~10 min) ──────────────────────
\section{Motivation and Background}

\begin{frame}{The Big Picture}
  \begin{itemize}
    \item Broad context — why this field matters
    \item The specific problem we address
    \item State of the art and its limitations
  \end{itemize}
  \note{For seminars, assume mixed expertise. Spend more time on context than in a conference talk.}
\end{frame}

\begin{frame}{Literature Context}
  \begin{itemize}
    \item Foundational work: [Author et al., Year]
    \item Recent advances: [Author et al., Year]
    \item Key gap this work fills
  \end{itemize}
  \note{Cite 5–8 key papers. Show you know the field deeply.}
\end{frame}

\begin{frame}{Research Questions}
  \begin{enumerate}
    \item Primary question / hypothesis
    \item Secondary question (if applicable)
    \item Expected contribution
  \end{enumerate}
  \note{Make the questions crisp and testable. Audience should remember these.}
\end{frame}

% ── Methods (5–7 slides, ~12 min) ────────────────────────────────────────
\section{Methods}

\begin{frame}{Study Design Rationale}
  \begin{itemize}
    \item Why this approach over alternatives
    \item Design trade-offs and their justification
    \item Connection to research questions
  \end{itemize}
  \note{Seminars allow deeper methods discussion. Explain your reasoning, not just the procedure.}
\end{frame}

\begin{frame}{Data and Participants}
  \begin{columns}[T]
    \begin{column}{0.48\textwidth}
      \textbf{Data Sources}
      \begin{itemize}
        \item Source 1: description, coverage
        \item Source 2: description, coverage
        \item Quality control procedures
      \end{itemize}
    \end{column}
    \begin{column}{0.48\textwidth}
      \framebox[\textwidth]{\parbox{0.9\textwidth}{\centering\vspace{2cm}Data Overview Figure\vspace{2cm}}}
    \end{column}
  \end{columns}
\end{frame}

\begin{frame}{Analysis Procedures}
  \begin{enumerate}
    \item Pre-processing pipeline
    \item Core analysis method
    \item Validation strategy
  \end{enumerate}
  \note{Walk through the pipeline step by step. Seminars reward methodological depth.}
\end{frame}

\begin{frame}{Analysis Plan}
  \begin{table}
    \caption{Analysis components mapped to research questions}
    \begin{tabular}{lll}
      \toprule
      Research Question & Analysis Method & Validation \\
      \midrule
      RQ1 & Method A & Cross-validation \\
      RQ2 & Method B & Sensitivity test \\
      \bottomrule
    \end{tabular}
  \end{table}
\end{frame}

% ── Results: Aim 1 (4–5 slides, ~8 min) ──────────────────────────────────
\section{Results: Aim 1}

\begin{frame}{Aim 1 — Overview}
  \begin{itemize}
    \item What we tested / explored
    \item Summary of findings
  \end{itemize}
\end{frame}

\begin{frame}{Aim 1 — Main Result}
  \begin{center}
    \framebox[0.8\textwidth]{\parbox{0.7\textwidth}{\centering\vspace{3cm}Aim 1 Key Figure\vspace{3cm}}}
  \end{center}
  \note{Let the figure speak. One slide, one message.}
\end{frame}

\begin{frame}{Aim 1 — Supporting Details}
  \begin{columns}[T]
    \begin{column}{0.48\textwidth}
      \framebox[\textwidth]{\parbox{0.9\textwidth}{\centering\vspace{2cm}Supporting Fig A\vspace{2cm}}}
    \end{column}
    \begin{column}{0.48\textwidth}
      \begin{itemize}
        \item Interpretation point 1
        \item Interpretation point 2
        \item Connection to hypothesis
      \end{itemize}
    \end{column}
  \end{columns}
\end{frame}

% ── Results: Aim 2 (4–5 slides, ~8 min) ──────────────────────────────────
\section{Results: Aim 2}

\begin{frame}{Aim 2 — Overview}
  \begin{itemize}
    \item What we tested / explored
    \item Summary of findings
  \end{itemize}
\end{frame}

\begin{frame}{Aim 2 — Main Result}
  \begin{center}
    \framebox[0.8\textwidth]{\parbox{0.7\textwidth}{\centering\vspace{3cm}Aim 2 Key Figure\vspace{3cm}}}
  \end{center}
\end{frame}

% ── Results: Aim 3 (4–5 slides, ~8 min) ──────────────────────────────────
\section{Results: Aim 3}

\begin{frame}{Aim 3 — Overview}
  \begin{itemize}
    \item What we tested / explored
    \item Summary of findings
  \end{itemize}
\end{frame}

\begin{frame}{Aim 3 — Main Result}
  \begin{center}
    \framebox[0.8\textwidth]{\parbox{0.7\textwidth}{\centering\vspace{3cm}Aim 3 Key Figure\vspace{3cm}}}
  \end{center}
\end{frame}

% ── Discussion (3 slides, ~7 min) ─────────────────────────────────────────
\section{Discussion}

\begin{frame}{Synthesis}
  \begin{itemize}
    \item How the three aims connect
    \item What the combined results tell us
    \item Comparison with prior work
  \end{itemize}
  \note{This is the intellectual core of the seminar. Connect the dots for the audience.}
\end{frame}

\begin{frame}{Implications}
  \begin{itemize}
    \item Theoretical implications
    \item Practical applications
    \item Policy relevance (if applicable)
  \end{itemize}
\end{frame}

\begin{frame}{Limitations and Caveats}
  \begin{itemize}
    \item Known limitation 1 and its impact
    \item Known limitation 2 and how we mitigate
    \item What this work does NOT claim
  \end{itemize}
  \note{Proactively addressing limitations builds credibility.}
\end{frame}

% ── Future Directions (1–2 slides) ────────────────────────────────────────
\section{Future Directions}

\begin{frame}{Next Steps}
  \begin{enumerate}
    \item Short-term: immediate follow-up study
    \item Medium-term: extending the framework
    \item Long-term: broader applications
  \end{enumerate}
  \note{Show you have a research trajectory. Good for seminar audiences evaluating your program.}
\end{frame}

% ── Conclusion ────────────────────────────────────────────────────────────
\section{Conclusion}

\begin{frame}{Summary}
  \begin{enumerate}
    \item Key finding 1 (from Aim 1)
    \item Key finding 2 (from Aim 2)
    \item Key finding 3 (from Aim 3)
  \end{enumerate}
  \vfill
  \small
  Code: \url{https://github.com/your/repo} \\
  Preprint: \url{https://arxiv.org/abs/xxxx.xxxxx}
\end{frame}

\begin{frame}{Acknowledgments}
  \small
  This work was supported by [Funding Agency, Grant \#].

  We thank [collaborators] for contributions and [reviewers] for feedback.

  % \printbibliography
\end{frame}

\begin{frame}[standout]
  Questions?
\end{frame}

% ══════════════════════════════════════════════════════════════════════════
\appendix

\begin{frame}{Backup: Full Regression Table}
  % Complete statistical results
\end{frame}

\begin{frame}{Backup: Model Fit Diagnostics}
  % Residual plots, convergence, fit statistics
\end{frame}

\begin{frame}{Backup: Additional Analyses}
  % Sensitivity analyses, robustness checks
\end{frame}

\begin{frame}{Backup: Code Example}
  \begin{lstlisting}[language=Python]
# Core analysis snippet
import numpy as np
results = model.fit(X_train, y_train)
print(f"R2 = {results.rsquared:.3f}")
  \end{lstlisting}
\end{frame}

\begin{frame}{Backup: Alternative Specifications}
  % Results under different assumptions or parameters
\end{frame}

\end{document}
